\documentclass[10pt,a4paper]{amsart}
\usepackage[margin=19mm]{geometry}
\usepackage{amsmath,amssymb,mathtools}
\usepackage[scr=rsfs,cal=euler]{mathalfa}
\usepackage{xfrac}
\usepackage[unicode,hidelinks,bookmarks=false]{hyperref}
\newcommand{\ie}{i.e.\ }
\newcommand{\cf}{cf.\ }
\newcommand{\resp}{respectively\ }
\newcommand{\Subsection}{\subsection}
\providecommand{\nobreakdash}{\nobreak\hbox{-}}
\setlength{\emergencystretch}{2em}
\multlinegap0pt
\usepackage{amssymb}
\usepackage{mathtools}
\usepackage[shortlabels]{enumitem}
\usepackage{tikz}
\newtheorem{lemma}{Lemma}
\newcommand{\D}{\mathcal D}
\renewcommand{\P}{\mathcal P}
\newcommand{\Z}{\mathbb Z}
\newcommand{\la}{\lambda}
\newcommand{\vecx}{\vec{x}}
\newcommand{\zeroes}{\vec{x}=(0^N)}
\title{Law of Large Numbers for continuous $N$-particle ensembles at fixed temperature}
\author{Cesar Cuenca, Jiaming Xu}
\date{Source: arXiv:2512.04394v2 (CC BY 4.0)}
\begin{document}
\maketitle

\noindent\emph{Source and licence.} Law of Large Numbers for continuous $N$-particle ensembles at fixed temperature. Version arXiv:2512.04394v2 (CC BY 4.0), distributed by arXiv under the Creative Commons Attribution 4.0 International licence, http://creativecommons.org/licenses/by/4.0/, as recorded in the arXiv licence metadata for this exact version and shown on the abstract page https://arxiv.org/abs/2512.04394v2. Full licence notice: \texttt{licenses/arxiv-2512.04394-CC-BY-4.0.txt}.\par\smallskip

\noindent\textit{Editorial scope.} Counted unit: the article's lemma lem:tech\_2 (source0.tex lines 517--526). The article's complete original proof of it is delivered with it (source0.tex lines 527--567, one \texttt{proof} environment of 41 lines). No mathematics is written, repaired or re-derived: the statement and the proof are the article's own lines, in the article's own notation, and no retained line is substituted or re-flowed.  No further source-local context is needed: every cross-reference of every retained line resolves inside the two delivered spans.  Nothing was omitted from any retained span, so the package carries no omission record.  No retained line contains a citation, so the wrapper carries no bibliography and no bibliography entry.  No retained line contains a figure environment, an image-inclusion command or drawing code, so no image is delivered and the package declares no asset dependency.  Editorial formatting only, in addition: an AMS \texttt{amsart} preamble that restates the article's own declarations and macros used by the retained lines, namely the environments \texttt{lemma} on separate counters, together with the standard packages those lines need.  The retained environments print in this document as Lemma 1 in order of appearance, whereas the article prints the same items with the numbering of its own full document.  The complete native source of the article as distributed in the arXiv e-print for this exact version is bundled unchanged as \texttt{sources/190327/source0.tex}.
\begin{lemma}\label{lem:tech_2}
    If for all partitions $\mu$, we have the limit
    \begin{equation}\label{eq:assumption_tech_2}
    \lim_{N\to\infty}{ N^{-|\mu|}\left[ \prod_{i=1}^{\ell(\mu)}{\D_i^{\mu_i}} \right] e^{F_{N,\theta}(\vecx)} \bigg|_{\vecx=(0^N)} } = \prod_{i=1}^{\ell(\mu)}{m_{\mu_i}},
    \end{equation}
    then for all partitions $\la$, we have
    \begin{equation}\label{eq:conclusion_tech_2}
    \lim_{N\to\infty} N^{-|\la|-\ell(\la)}\left[ \prod_{i=1}^{\ell(\la)}{\P_{\la_i}} \right] e^{F_{N,\theta}(\vecx)} \bigg|_{\vecx=(0^N)} = \prod_{i=1}^{\ell(\la)}{m_{\la_i}}.
    \end{equation}
\end{lemma}

\begin{proof}
Assume that Eqn.~\eqref{eq:assumption_tech_2} holds for all $\mu$. Let us take any partition $\la$ and prove Eqn.~\eqref{eq:conclusion_tech_2}.

If $\ell(\la)=1$, then $\la=(n)$, for some $n\in\Z_{\ge 1}$.
Note that $\D_i^n e^{F_{N,\theta}(\vecx)}$ differs from $\D_j^n e^{F_{N,\theta}(\vecx)}$ only by the change of variables $x_i\leftrightarrow x_j$ because of the symmetry of $F_{N,\theta}(\vecx)$; in particular,
\begin{equation}\label{eq:equality_1}
\D_i^n e^{F_{N,\theta}(\vecx)}\big|_{\vecx=(0^N)} = \D_j^n e^{F_{N,\theta}(\vecx)}\big|_{\vecx=(0^N)}.
\end{equation}
As a result,
\begin{equation*}
N^{-n-1}\,\P_n e^{F_{N,\theta}(\vecx)} \Big|_{\vecx=(0^N)} = N^{-n-1}\sum_{i=1}^n{\D_i^n} e^{F_{N,\theta}(\vecx)} \Big|_{\vecx=(0^N)} = N^{-n}\,\D_1^n e^{F_{N,\theta}(\vecx)}\Big|_{\vecx=(0^N)},
\end{equation*}
proving that the desired Eqn.~\eqref{eq:conclusion_tech_2} follows at once from Eqn.~\eqref{eq:assumption_tech_2} for $\mu=(n)$.

\smallskip

If $\ell(\la)=2$, then $\la=(n_1,n_2)$, for some positive integers $n_1\ge n_2$. As above, we deduce that $\D_i^{n_2}\D_j^{n_1} e^{F_{N,\theta}(\vecx)}\big|_{\vecx=(0^N)} = \D_2^{n_2}\D_1^{n_1} e^{F_{N,\theta}(\vecx)}\big|_{\vecx=(0^N)}$, whenever $i\ne j$.
Together with~\eqref{eq:equality_1}, we obtain
\begin{multline*}
N^{-n_1-n_2-2}\,\P_{n_2}\P_{n_1} e^{F_{N,\theta}(\vecx)} = N^{-n_1-n_2-2}\sum_{i=1}^N{\D_i^{n_2}}\sum_{j=1}^N{\D_j^{n_1}} e^{F_{N,\theta}(\vecx)}\\
= N^{-n_1-n_2-2}\cdot N(N-1)\D_2^{n_2}\D_1^{n_1} e^{F_{N,\theta}(\vecx)} + N^{-n_1-n_2-2}\cdot N\,\D_1^{n_1+n_2} e^{F_{N,\theta}(\vecx)}\\
= \left(1-\frac{1}{N}\right)\cdot N^{-n_1-n_2}\,\D_2^{n_2}\D_1^{n_1}e^{F_{N,\theta}(\vecx)} + \frac{1}{N}\cdot N^{-n_1-n_2}\,\D_1^{n_1+n_2}e^{F_{N,\theta}(\vecx)}.
\end{multline*}
By assumption~\eqref{eq:assumption_tech_2}, the first term above converges to $m_1m_2$ and the second one to $0$, as $N\to\infty$.
Then Eqn.~\eqref{eq:conclusion_tech_2} follows.

\smallskip

If $\ell(\la)=\ell>2$, the proof is similar.
Indeed,
\begin{equation}\label{eq:equality_2}
\prod_{i=1}^{\ell(\la)}{\P_{\la_i}} = \prod_{i=1}^{\ell(\la)}{\sum_{j=1}^N{\D_j^{\la_i}}} = \sum_{\substack{1\le j_1,\dots,j_{\ell(\la)}\le N \\ j_1,\dots,j_{\ell(\la)}\text{ are distinct}}}{\D_{j_1}^{\la_1}\cdots\D_{j_{\ell(\la)}}^{\la_{\ell(\la)}}} + p(\D_1,\dots,\D_N),
\end{equation}
where $p(x_1,\dots,x_N)$ is a symmetric polynomial, where each monomial has at most $\ell(\la)-1$ distinct variables $x_j$.
Again, we have that if $j_1,\dots,j_{\ell(\la)}$ are distinct, then $\D_{j_1}^{\la_1}\cdots\D_{j_{\ell(\la)}}^{\la_{\ell(\la)}} e^{F_{N,\theta}(\vecx)}\big|_{\vecx=(0^N)}$ is equal to $\D_1^{\la_1}\cdots\D_{\ell(\la)}^{\la_{\ell(\la)}} e^{F_{N,\theta}(\vecx)}\big|_{\vecx=(0^N)}$. This shows, as in the case when $\ell(\la)=2$, that~\eqref{eq:equality_2} gives
\begin{multline*}
N^{-|\la|-\ell(\la)} \left[\prod_{i=1}^{\ell(\la)}{\P_{\la_i}}\right]  e^{F_{N,\theta}(\vecx)} \bigg|_{\vecx=(0^N)} = \left( 1 + O\Big( \frac{1}{N} \Big) \right)\cdot N^{-|\la|} \D_1^{\la_1}\cdots\D_{\ell(\la)}^{\la_{\ell(\la)}} e^{F_{N,\theta}(\vecx)}\big|_{\zeroes}\\
+ \sum_{j=1}^{\ell(\la)-1} \sum_{\substack{k_1,\dots,k_j\in\Z_{\ge 1} \\ k_1+\dots+k_j=|\la|}} O\Big( \frac{1}{N} \Big)\cdot N^{-k_1-\,\dots\,-k_j} \D_1^{k_1}\cdots\D_j^{k_j} e^{F_{N,\theta}(\vecx)}\big|_{\zeroes},
\end{multline*}
and then assumption~\eqref{eq:assumption_tech_2} gives the desired equation~\eqref{eq:conclusion_tech_2}.
\end{proof}

\end{document}
